\chapter{Introduction générale}
\label{chap:introduction}

Dans un monde professionnel de plus en plus compétitif, la qualité et la pertinence d'un curriculum vitae (CV) jouent un rôle déterminant dans le succès d'une candidature. Chaque offre d'emploi possède ses propres exigences, et un CV générique a souvent moins d'impact qu'un CV adapté spécifiquement au poste visé. Cependant, personnaliser manuellement son CV pour chaque candidature est un processus chronophage et fastidieux.

C'est dans ce contexte que s'inscrit le projet \projectname{}, une plateforme web intelligente de génération de CV assistée par intelligence artificielle. L'objectif est de permettre aux candidats de générer un CV personnalisé et optimisé pour une offre d'emploi donnée en quelques secondes, tout en assurant une compatibilité optimale avec les systèmes de suivi des candidatures (ATS).

Ce rapport présente l'ensemble du travail réalisé par notre équipe dans le cadre du projet d'intégration de la 1ère année Génie Informatique, sous la supervision de \professeur.

Le rapport est structuré en trois chapitres principaux. Le premier chapitre pose le contexte général du projet, présente les objectifs, la méthodologie de gestion adoptée ainsi que l'équipe de réalisation. Le deuxième chapitre est consacré à l'analyse et la conception du système, incluant l'étude des besoins, l'identification des acteurs et la modélisation. Enfin, le troisième chapitre détaille la réalisation technique du projet, couvrant l'architecture, les outils utilisés, le cycle DevSecOps, l'implémentation et le déploiement.
